% !PRINTED 623
\weberpage{623}

%<<CONT
in a similar way one shows that no two others of these
subsidiary groups are equal to one another.

According to this, the values of $u$ conjugate within $A$ are
obtained by applying the permutations

\[
(1,\,2)(3,\,4), \quad t(0,1) = (0,1,2,4,3), \quad t(0,2) = (0,2,4,3,1)
\]
\[
t(0,3) = (0,3,1,2,4), \quad t(0,4) = (0,4,3,1,2)\text{:}
\]

\weq{11}{
\begin{aligned}
u_1 &= x_0x_1 + x_1x_2 + x_2x_3 + x_3x_4 + x_4x_0, \\
u_2 &= x_0x_2 + x_2x_1 + x_1x_4 + x_4x_3 + x_3x_0, \\
u_3 &= x_1x_2 + x_2x_4 + x_4x_0 + x_0x_3 + x_3x_1, \\
u_4 &= x_2x_0 + x_0x_4 + x_4x_1 + x_1x_3 + x_3x_2, \\
u_5 &= x_3x_2 + x_2x_4 + x_4x_1 + x_1x_0 + x_0x_3, \\
u_6 &= x_4x_2 + x_2x_0 + x_0x_1 + x_1x_3 + x_3x_4,
\end{aligned}
}

and the corresponding $u_1', u_2', u_3', u_4', u_5', u_6'$ are found by
applying the permutation $t$ to the $u_1$, and then proceeding with $u_1'$
just as was done with $u_1$, or also simply by uniting the pairs
missing from each $u$ into the corresponding $u'$,
so that for all these $u$ the condition (7) is satisfied.

\weq{12}{
\begin{aligned}
u_1' &= x_0x_2 + x_0x_3 + x_1x_3 + x_1x_4 + x_2x_4, \\
u_2' &= x_0x_1 + x_0x_4 + x_1x_3 + x_2x_3 + x_2x_4, \\
u_3' &= x_0x_1 + x_0x_2 + x_2x_3 + x_3x_4 + x_1x_4, \\
u_4' &= x_0x_1 + x_0x_3 + x_1x_2 + x_2x_4 + x_3x_4, \\
u_5' &= x_0x_2 + x_0x_4 + x_1x_2 + x_1x_3 + x_3x_4, \\
u_6' &= x_0x_3 + x_0x_4 + x_1x_2 + x_1x_4 + x_2x_3.
\end{aligned}
}

The six quantities (8): $y_1, y_2, y_3, y_4, y_5, y_6$ are the roots of an
equation of the sixth degree, whose coefficients depend rationally on $a, b, c, d, e$
and $\sqrt{\varDelta}$. Since the $y$ all change their sign when
$\sqrt{\varDelta}$ changes its sign, this equation has the form

\weq{12}{
y^6 + a_2y^4 + a_4y^2 + a_6 - \sqrt{\varDelta}(a_1y^5 + a_3y^3 + a_5y) = 0, % [sic] duplicate equation number (12); printed page repeats (12) rather than continuing to (13)
}

in which the $a$ are rational functions of the coefficients of $f(x)$.
But they must also be integral functions of these coefficients;
for $a_2, a_4, \alpha_6, a_1\sqrt{\varDelta}, a_3\sqrt{\varDelta}, a_5\sqrt{\varDelta}$ are integral functions of the % [sic] printed as $\alpha_6$ (lowercase alpha) rather than $a_6$
$x$, which can indeed be regarded here as independent variables, and
therefore the three last must be divisible by the difference product $\sqrt{\varDelta}$,
and then $a_2, a_4, a_6, a_1, a_3, a_5$ must prove to be integral
symmetric functions of the $x$. If one determines the
degrees with respect to the $x$, then, since the $y$ are of the second
degree, there result for
%CONT>>
